% ── §6 Formal Verification in Lean 4 ──
\section{Formal Verification in Lean 4}
\label{sec:theory-lean}

\subsection{Motivation}
The encoding layer through which biological sequences enter the statistical
machinery of coevolution analysis involves several non-obvious design choices:
the assignment of integer codes to amino acids, the construction of Gray-code
transformations, the definition of Hamming distances on encoded values, and
the semantics of implicant matching on Karnaugh-map cells.  Errors in any of
these choices propagate silently through the pipeline and can produce
plausible-looking but incorrect results.  To address this, we formalise all
load-bearing definitions and properties in the Lean~4 theorem
prover~\cite{lean4}, establishing machine-checkable proofs that the
implementation satisfies its specification.

\subsection{Verification Architecture}
Our formalisation comprises 118 theorems distributed across five modules in
two repositories: \texttt{lean\_proofs/} (106 theorems covering Gray-code
fundamentals, $k$-mer indexing, amino-acid encodings, contact-map topology,
and encoding equivalence) and \texttt{ContactCircuits.lean} (12 theorems
covering separation-bin monotonicity, level-1 cell-index boundedness and
injectivity, implicant-cover completeness and soundness, and padding safety).
All proofs are constructed without \texttt{sorry} and rely only on the
standard axioms \texttt{propext}, \texttt{Quot.sound}, and (where
\texttt{native\_decide} is used) the trusted-evaluation axiom inherent to
bounded-domain computation.

\subsection{Key Theorems}

\subsubsection{Gray-Code Properties}
The following theorems establish the foundational properties of the reflected
binary Gray code over fixed-width domains:
\begin{itemize}[nosep]
	\item \texttt{cc\_gray\_consecutive\_adj}: consecutive integers differ by
	      exactly one bit under Gray transformation.
	\item \texttt{cc\_hd\_symm}: the Gray-Hamming distance is symmetric.
	\item \texttt{cc\_he\_dist1\_ordered\_eq\_80}: exactly 80 ordered pairs of
	      He-coded residues lie at Gray distance~1.
	\item \texttt{cc\_he\_dist\_counts}: the full ordered distribution at
	      distances 0--5 equals $[20, 80, 132, 112, 48, 8]$.
\end{itemize}

\subsubsection{Implicant Semantics}
The following theorems formalise the semantics of Quine--McCluskey output,
matching the runtime assertions made by every circuit-extraction script:
\begin{itemize}[nosep]
	\item \texttt{cc\_cover\_complete}: every on-set cell is matched by some
	      prime implicant.
	\item \texttt{cc\_off\_avoiding}: no prime implicant matches any off-set
	      cell.
\end{itemize}

\subsubsection{Padding Safety}
The 32$\times$32 padding of a 20$\times$20 map introduces don't-care border
cells whose codes exceed the canonical residue range.  The following theorems
guarantee that these border cells cannot leak into decoded rules:
\begin{itemize}[nosep]
	\item \texttt{cc\_fixed\_match\_unique}: a fully-fixed 10-bit implicant
	      matches exactly one cell, namely its own packed index.
	\item \texttt{cc\_padding\_safety}: an implicant decoding to canonical
	      residues ($< 20$ on both axes) fires only inside the $20\times20$
	      block.
\end{itemize}

\subsection{Axiom Audit}
We verify the axiom dependency graph for all 12 theorems in the
\texttt{ContactCircuits} module using Lean's \texttt{\#print axioms} command.
Every theorem depends only on \texttt{propext}, \texttt{Quot.sound}, and
(where applicable) the native-decide evaluation axiom; no theorem depends on
\texttt{sorryAx}.  This matches the trust profile of the existing 106 theorems
in the companion repository.
